\documentclass[10pt,a4paper]{article}
\usepackage[margin=2.2cm]{geometry}
\usepackage{fontspec}
\usepackage{kotex}
\setmainfont{Apple SD Gothic Neo}
\setsansfont{Apple SD Gothic Neo}
\usepackage{amsmath,amssymb}
\usepackage{graphicx}
\usepackage{booktabs}
\usepackage{caption}
\usepackage[dvipsnames]{xcolor}
\usepackage[colorlinks=true,linkcolor=MidnightBlue,citecolor=MidnightBlue,
            urlcolor=MidnightBlue]{hyperref}
\usepackage{titlesec}
\titleformat{\section}{\large\bfseries}{\thesection}{0.6em}{}
\captionsetup{font=small,labelfont=bf}
\setlength{\parskip}{0.35em}
\setlength{\parindent}{0pt}

\usepackage{mdframed}
\newmdenv[linewidth=0.4pt,linecolor=gray!60,backgroundcolor=gray!5,
          innertopmargin=6pt,innerbottommargin=6pt,skipabove=8pt,skipbelow=8pt]{claimbox}
\newcommand{\claim}[2]{\begin{claimbox}\textbf{주장.} #1\par\smallskip
  \textcolor{gray!70}{\textbf{반증 조건.} #2}\end{claimbox}}

\title{\bfseries 간발의 불확인\\[2pt]
\large $0.0014$ 가 모자랐다 --- 그리고 병기한 부트스트랩이 자를 또 부러뜨렸다}
\author{월드모델 실험실 · 노트 $826$}
\date{2026-08-07}
\begin{document}\maketitle

\section{한 줄}

티처 \#$2$ 의 ``게이트 도전 실험 $0$건 --- 회피의 통계적 서명'' 지적을 받아
$20$노트 만의 채택-가능 실험을 발동했다: T$1$ 미결에 등록돼 있던 \textbf{아이돌
라우팅 확인 사이클}. 결과는 \textbf{간발의 불확인}이다 --- 챔피언 씨앗 $1{\sim}4$
의 $\Delta$(합동$-$릿지)가 \textbf{$4/4$ 음수}인데 평균 $-0.0720$ 이 등록 문턱
$-0.0734$ 를 $0.0014$ 차로 못 넘었다. 등록 규칙 그대로: 합동 유지 · 아이돌
재측정 금지 $3$사이클.

\begin{figure}[t]\centering
\includegraphics[width=0.97\textwidth]{figs/hair.pdf}
\caption{왼쪽: 씨앗별 $\Delta$ 와 문턱. 오른쪽: 상수 $2\sigma$ 띠 대 짝
부트스트랩 $95\%$ --- 폭이 $3$배이고 $0$ 을 문다.}
\end{figure}

\section{병기가 본론을 삼켰다}

\claim{\textbf{상수 자($2\sigma = 2{\times}0.262/\sqrt{n}$)가 실측으로
반증됐다.} 사전등록이 병기로 걸어 둔 짝 행 부트스트랩 $1{,}000$회의 $\Delta$
$95\%$ 구간은 $[-0.293, +0.137]$ --- 폭이 상수 띠($\pm 0.0734$)의 약 $3$배이고
$0$ 을 여유 있게 문다. 즉 논문 $182$ 의 ``아이돌 릿지승($-0.086$)''은
\textbf{부러진 자 위의 판정}이었다. $n{=}51$ 순위상관 차이의 분산을 상수가
크게 과소평가한다 --- 순위상관은 정규 근사보다 무겁고 두 모형의 짝 상관도
무시된다. 티처 \#$2$ 가 지목한 약점 $1$ 의 즉석 실측 확인이다.}
{두꺼운 도메인의 살찐 마진(도서 $7\times$ · 펀딩 $10\times$ · 애니 $7\times$)은
살아남을 것이다. 다음 사전등록: \textbf{짝 부트스트랩으로 $822$(b) $10$도메인
전부 재채점} --- 예상은 아이돌 릿지승 $\to$ 동점 강등 · 합동승 $5$ 유지이고,
그러면 지도의 문장이 \emph{``합동은 어디서도 유의하게 해롭지 않고 $5$곳에서
유의하게 이롭다''} 로 오히려 단순해지고 강해진다. 규약 후보(재채점 후):
모델 대 모델 짝 비교에 상수 자 금지.}

\section{부록 --- 티처의 값과 한계}

티처 \#$2$ 의 유효 지적 넷을 반영했다($825$ 매처 조건부 · 사후 서사 경고 ·
상수 자 · 게이트 도전). 동시에 \textbf{티처의 한계도 실측}됐다: 제안 B(웹툰
후보 열)는 노트 $777/787/799$ 가 이미 닫은 갈래였다 --- 프레시 눈은 아카이브
깊이가 없으므로, \textbf{티처 제안은 대장 교차검증 후 채택}이 절차가 된다.
곁 배선: 유튜브 전향 폴러(\texttt{ingest/yt\_poll.py} · RSS$+$워치 · 키 불요)
완성 --- 남은 것은 사람 게이트(대상 목록)뿐. 판 $\rho$ $0.4689$ 그대로.

\end{document}
